\section{Training the CNN: Loss, Gradient Descent \& Generalisation}
\label{sec:training}

The shipped Cadence monitor decides freeze episodes with a transparent freeze-index
rule. The CNN of the previous section --- \emph{FoGNet}, a 1-D convolutional network
of roughly \num{18000} parameters --- is the ``given more time'' upgrade: instead of
hand-reading one ratio, it \emph{learns} from data what a freeze window looks like.
But a network with eighteen thousand free numbers does not know anything at birth; it
has to be \emph{taught}. This section is the story of that teaching. We need three
ingredients: a way to turn the network's raw opinion into a probability (the
\textbf{sigmoid}), a way to score how wrong that probability is (the \textbf{loss}),
and a way to nudge the eighteen thousand numbers so the score improves (\textbf{gradient
descent}). Then comes the part that actually decides whether the model is trustworthy:
\textbf{how we evaluate it} without fooling ourselves.

\begin{intuition}
Training is just \emph{practising with an answer key}. We show FoGNet thousands of
4-second accelerometer windows from real Parkinson's patients, each stamped
``freeze'' or ``calm'' by a clinician. For every window the network guesses, we measure
how badly it guessed, and we adjust its internal knobs a hair in the direction that
would have made the guess better. Repeat a few hundred thousand times and the knobs
settle into a configuration that recognises freezing it has never seen before.
\end{intuition}

\subsection{From a raw score to a probability: the sigmoid}

FoGNet's final layer emits two raw numbers --- one \emph{logit} arguing ``freeze'',
one arguing ``calm''. Their difference $z = z_{\text{freeze}} - z_{\text{calm}}$ is a
single signed score: positive leans freeze, negative leans calm, zero is a coin-flip.
A score is not yet a probability, though --- $z$ can be any real number, while a
probability must live in $[0,1]$. The \textbf{sigmoid} (equivalently, the 2-class
\emph{softmax}) is the squashing function that maps the whole real line onto that
interval:
\begin{equation}
\hat{p} \;=\; P(\text{freeze}\mid \text{window}) \;=\; \sigma(z) \;=\; \frac{1}{1+e^{-z}}.
\label{eq:sigmoid}
\end{equation}
A large positive logit gives $\hat p \to 1$; a large negative one gives $\hat p \to 0$;
and $z=0$ sits exactly at $\hat p = 0.5$. The final ``freeze / calm'' call is made by
comparing $\hat p$ against a decision threshold $t$ --- and, crucially, $t$ need
\emph{not} be $0.5$. Because a missed freeze is far costlier than a false alarm, we are
free to slide $t$ downward to favour sensitivity, exactly as the shipped FI rule biases
its trip line. Choosing $t$ honestly is the subject of the calibration discussion at
the end of this section.

Figure~\ref{fig:sigmoid_bce} (left) shows the curve. Two features matter for what
follows. First, the steep middle: near $z=0$ a small change in the score moves the
probability a lot, so the network can be \emph{decisive}. Second, the flat tails: once
$z$ is far from zero the curve \emph{saturates} and its slope approaches zero. A flat
slope means a near-zero gradient, and a near-zero gradient means almost no learning
signal --- a fact that will return when we discuss how training can stall.

\begin{figure}[t]
  \centering
  \begin{subfigure}[t]{0.49\textwidth}
    \centering
    \includegraphics[width=\linewidth]{sigmoid}
    \caption{The sigmoid maps any score $z$ to a probability in $[0,1]$; it is steep
    in the middle (decisive) and flat in the tails (saturated, gradient $\approx 0$).}
    \label{fig:sigmoid}
  \end{subfigure}\hfill
  \begin{subfigure}[t]{0.49\textwidth}
    \centering
    \includegraphics[width=\linewidth]{bce_loss}
    \caption{Binary cross-entropy. Being confidently \emph{wrong} ($\hat p$ near the
    opposite label) sends the loss towards infinity; being confidently \emph{right}
    costs almost nothing.}
    \label{fig:bce}
  \end{subfigure}
  \caption{The two functions that turn ``guessing'' into ``learning''. The sigmoid
  (left) produces a probability; binary cross-entropy (right) prices how bad that
  probability is given the true label.}
  \label{fig:sigmoid_bce}
\end{figure}

\subsection{The loss: punishing confident mistakes}

We now need a number that says \emph{how wrong} a probability $\hat p$ is, given the
true label $y \in \{0,1\}$ ($y=1$ for a real freeze window, $y=0$ for calm gait).
FoGNet is trained with \textbf{binary cross-entropy} (BCE), also called log-loss:
\begin{equation}
\mathcal{L}(\hat p, y) \;=\; -\Big[\, y\,\log \hat p \;+\; (1-y)\,\log(1-\hat p)\,\Big].
\label{eq:bce}
\end{equation}

\begin{mathsbox}
Equation~\eqref{eq:bce} looks like two terms, but for any single window only one of
them is ever ``on'', because $y$ is either $0$ or $1$.
\begin{itemize}
  \item \textbf{A true freeze ($y=1$).} The second term vanishes and the loss collapses
  to $\mathcal{L} = -\log \hat p$. If the network correctly said $\hat p = 0.99$, the
  loss is $-\log(0.99)\approx 0.01$ --- almost free. If it confidently insisted
  ``calm'' with $\hat p = 0.01$, the loss is $-\log(0.01)\approx 4.6$ --- a huge
  penalty.
  \item \textbf{A calm window ($y=0$).} Now the first term vanishes and
  $\mathcal{L} = -\log(1-\hat p)$, which is the mirror image: cheap when $\hat p$ is
  near $0$, explosive when the network falsely cries ``freeze''.
\end{itemize}
The $\log$ is what makes BCE the right tool. A plain squared-error would mildly scold a
confident blunder; the logarithm makes the penalty diverge to $+\infty$ as $\hat p$
heads to the \emph{wrong} certainty. The network therefore learns not just to be right,
but to be \emph{calibrated} --- to reserve its confidence for cases that deserve it.
\end{mathsbox}

Figure~\ref{fig:sigmoid_bce} (right) draws both branches. The training objective is the
\emph{average} of Equation~\eqref{eq:bce} over the whole batch of labelled windows; that
single scalar is what the optimiser tries to push downhill.

\begin{intuition}
BCE is a teacher who marks two things at once: \emph{was the answer right}, and
\emph{how sure were you}. A hesitant wrong answer is forgiven; a confident wrong answer
is hammered. That is precisely the behaviour we want from a fall-risk monitor --- a
model that shouts ``freeze!'' must have earned the right to shout.
\end{intuition}

\subsection{Gradient descent: rolling the weights downhill}

The loss is a function of the network's \num{18000} weights: change a weight, and the
average loss changes. Picture those weights as coordinates on a vast, hilly landscape
whose \emph{height} is the loss. Training is the search for the lowest valley.
\textbf{Backpropagation} (the chain rule applied layer by layer) computes the
\emph{gradient} $\nabla_{\!\boldsymbol\theta}\mathcal{L}$ --- the direction of steepest
\emph{uphill} at the current spot. To go down, we step the opposite way:
\begin{equation}
\boldsymbol\theta \;\leftarrow\; \boldsymbol\theta \;-\; \eta\,
\nabla_{\!\boldsymbol\theta}\mathcal{L},
\label{eq:sgd}
\end{equation}
where $\boldsymbol\theta$ collects all the weights and $\eta$ is the
\textbf{learning rate} --- the size of each step. (Cadence trains with the AdamW
optimiser, a refined cousin of Equation~\eqref{eq:sgd} that adapts the step per weight,
but the ``roll downhill'' picture is identical.) Repeat the read-the-slope /
take-a-step loop over many \emph{epochs} --- one epoch being a full pass over the
training windows --- and the weights descend towards a good valley.

\begin{figure}[t]
  \centering
  \includegraphics[width=0.74\textwidth]{gradient_descent}
  \caption{Gradient descent on a 2-D cartoon of the loss surface (the real surface is
  $\sim$\num{18000}-dimensional). The contours are loss; the star is the best weights.
  A learning rate too small \emph{crawls}; one too large \emph{zig-zags} and may never
  settle; a well-chosen rate \emph{converges} straight to the minimum.}
  \label{fig:gd}
\end{figure}

Figure~\ref{fig:gd} makes the role of the learning rate vivid. If $\eta$ is too small
the optimiser inches along and may run out of epochs before reaching the valley. If
$\eta$ is too large it overshoots, ricochets across the basin and can diverge. The
``good rate'' path walks almost straight to the star. This is also where the saturated
sigmoid tails of Figure~\ref{fig:sigmoid} bite: where the curve is flat the gradient is
tiny, the step in Equation~\eqref{eq:sgd} barely moves, and learning crawls --- one
reason normalisation layers (below) are used to keep activations in the lively middle
of the curve.

\begin{takeaway}
Loss says \emph{how wrong}; the gradient says \emph{which way is wrong}; the learning
rate says \emph{how big a correction to make}. Iterate, and eighteen thousand random
numbers turn into a freeze detector.
\end{takeaway}

\subsection{Over-fitting and the regularisers}

A network with \num{18000} parameters is powerful enough to \emph{memorise} its training
windows --- including their noise and the quirks of the particular patients it saw.
A memoriser scores beautifully on data it has already seen and fails on anyone new;
this is \textbf{over-fitting}. FoGNet leans on two built-in defences, visible in its
architecture:
\begin{itemize}
  \item \textbf{Dropout (rate 0.3)} in the classifier head. During training, a random
  $30\%$ of the units are switched off on each pass, so no single neuron can become a
  load-bearing memoriser; the network is forced to spread its evidence across many
  redundant features. At inference all units are active again.
  \item \textbf{Batch normalisation} after every convolution. BatchNorm re-centres and
  re-scales each layer's activations, which keeps signals in the responsive middle of
  the nonlinearity (avoiding the dead saturated tails), steadies the gradients, and
  lets training use a larger, faster learning rate.
\end{itemize}
The tell-tale sign of over-fitting is a \emph{gap}: the training loss keeps falling
while performance on held-out data stops improving --- or worsens. We therefore watch a
validation metric every epoch and keep the checkpoint that is best on data the network
was \emph{not} trained on, a practice called \textbf{early stopping}.

\begin{figure}[t]
  \centering
  \begin{subfigure}[t]{0.55\textwidth}
    \centering
    \includegraphics[width=\linewidth]{train_curve}
    \caption{Training curve. The cross-entropy loss (red) falls epoch after epoch while
    validation sensitivity and specificity climb. The starred epoch is the best
    validation checkpoint --- the weights that actually ship.}
    \label{fig:train}
  \end{subfigure}\hfill
  \begin{subfigure}[t]{0.43\textwidth}
    \centering
    \includegraphics[width=\linewidth]{calibration}
    \caption{Reliability diagram (glass-box model). Points on the dashed line mean a
    stated probability matches the observed freeze frequency; the histogram shows how
    many windows land in each confidence bin.}
    \label{fig:calib}
  \end{subfigure}
  \caption{Watching learning happen (left) and checking the probabilities are honest
  (right). Early stopping reads off the left panel; threshold choice reads off the
  right.}
  \label{fig:train_calib}
\end{figure}

Figure~\ref{fig:train_calib} (left) is FoGNet's actual training run on the Daphnet
ankle data: loss slides down as the validation sensitivity and specificity rise, and
the highlighted epoch marks the checkpoint with the best validation balance --- the one
that ships rather than the lowest-loss one.

\subsection{Evaluating without fooling ourselves: Leave-One-Subject-Out}

This is the most important methodological decision in the whole pipeline, and the
easiest to get catastrophically wrong. The obvious way to test a model is to shuffle all
the windows, hold back a random fraction, train on the rest, and report the score on the
held-back part. For freeze detection this is a \emph{trap}.

\begin{keyidea}
\textbf{A random window split leaks the person, not the freeze.} Consecutive 4-second
windows from one patient overlap and share that individual's gait fingerprint --- their
stride rhythm, their tremor, the exact way their sleeve sits on the ankle. Split them
randomly and almost every test window has a near-twin in the training set from the
\emph{same body}. The model can then ace the test by recognising \emph{who} is walking
rather than \emph{whether} they are freezing --- and that skill evaporates the moment a
new patient straps the device on. The honest question is never ``can it recognise this
patient again?'' but ``can it recognise freezing in someone it has \emph{never met}?''
\end{keyidea}

The fix is \textbf{Leave-One-Subject-Out (LOSO)} cross-validation. With $N$ patients we
run $N$ rounds; in each, one patient is held out entirely as the test set and the model
is trained on the other $N-1$. Every patient takes a turn as the unseen subject, and we
pool the $N$ held-out verdicts into one confusion matrix. No window from a test patient
ever touches training, so the score genuinely answers the generalisation question.
Figure~\ref{fig:loso} sketches the protocol.

\begin{figure}[t]
  \centering
  \resizebox{0.94\textwidth}{!}{%
  \begin{tikzpicture}[font=\sffamily,
      patient/.style={rectangle,rounded corners=2pt,minimum width=1.05cm,
        minimum height=0.62cm,draw=rule,line width=0.5pt},
      train/.style={patient,fill=palenavy,text=ink},
      test/.style={patient,fill=palefreeze,draw=freeze,text=freeze,line width=0.8pt},
      rlabel/.style={font=\normalsize\bfseries,text=navy,anchor=east}]

    \def\np{8}      % patients per row
    \def\dx{1.22}   % horizontal spacing

    % three illustrative rounds + an ellipsis row
    \foreach \round/\hold/\yy in {1/1/0, 2/2/-0.95, 3/3/-1.90} {
      \node[rlabel] at (-0.55,\yy) {Round \round};
      \foreach \p in {1,...,\np} {
        \pgfmathtruncatemacro{\istest}{ifthenelse(\p==\hold,1,0)}
        \ifnum\istest=1
          \node[test] (n\round-\p) at ({\p*\dx},\yy) {\footnotesize P\p};
        \else
          \node[train] (n\round-\p) at ({\p*\dx},\yy) {\footnotesize P\p};
        \fi
      }
    }
    % ellipsis row
    \node[rlabel] at (-0.55,-2.85) {$\vdots$};
    \foreach \p in {1,...,\np} {
      \node[font=\large,text=mute] at ({\p*\dx},-2.78) {$\vdots$};
    }
    % last round: hold out P8
    \node[rlabel] at (-0.55,-3.75) {Round $N$};
    \foreach \p in {1,...,\np} {
      \pgfmathtruncatemacro{\istest}{ifthenelse(\p==8,1,0)}
      \ifnum\istest=1
        \node[test] at ({\p*\dx},-3.75) {\footnotesize P\p};
      \else
        \node[train] at ({\p*\dx},-3.75) {\footnotesize P\p};
      \fi
    }

    % legend (stacked rows so long labels never collide; dark, legible labels —
    % the swatch boxes carry the blue/red colour key)
    \node[train,minimum width=0.8cm] (lg1) at (2.1,-4.95) {};
    \node[anchor=west,font=\normalsize\bfseries,text=ink] at (2.75,-4.95)
      {train on these $N{-}1$ patients};
    \node[test,minimum width=0.8cm] (lg2) at (2.1,-5.85) {};
    \node[anchor=west,font=\normalsize\bfseries,text=ink] at (2.75,-5.85)
      {held-out test patient (never seen)};
  \end{tikzpicture}}
  \caption{Leave-One-Subject-Out cross-validation. Each round holds out one whole
  patient (red) and trains on all the others (blue); rotating the held-out subject
  through everyone measures performance on bodies the network has never met. Pooling the
  held-out verdicts gives an unbiased confusion matrix.}
  \label{fig:loso}
\end{figure}

Run this way, FoGNet's pooled LOSO confusion over \num{8982} held-out windows gives
$\mathrm{TN}=\num{6842}$, $\mathrm{FP}=\num{1285}$, $\mathrm{FN}=248$,
$\mathrm{TP}=607$ --- a sensitivity of $0.71$, specificity of $0.84$ and accuracy of
$0.83$. These are honest, new-patient numbers, not the inflated figures a random split
would have printed. The interpretable 9-feature glass-box alternative, evaluated under
the same protocol, reaches sensitivity $0.75$ and specificity $0.76$ at its operating
point, with a threshold-free ROC-AUC of $0.82$.

\subsubsection*{Nested LOSO for the decision threshold}

There is one more leak to seal. The decision threshold $t$ of
Section~\ref{sec:training} is itself a parameter we tune --- and if we tune it by peeking
at the very test patients we then report on, we have quietly leaked the answer again.
\textbf{Nested LOSO} closes this: inside each outer round, the threshold is chosen on a
further held-out \emph{inner} fold drawn only from the training patients, never from the
outer test subject. The glass-box model's reported operating point ($t^\star \approx
0.10$) is selected this way, so even the threshold is honest. The headline numbers we
quote are therefore what a genuinely new wearer should expect, not a flattering
self-portrait.

\subsection{Calibration: do the probabilities mean what they say?}

A final, subtle question: when the model says ``$\hat p = 0.7$'', do roughly $70\%$ of
such windows actually turn out to be freezes? A model can rank windows well yet state
\emph{mis-scaled} probabilities --- which matters, because the Freeze-Margin meter on the
live dashboard and the choice of $t$ both assume the numbers are meaningful. We check
this with a \textbf{reliability diagram}: bin the predictions by stated probability and
plot, per bin, the observed freeze frequency against the predicted one.
Figure~\ref{fig:calib} shows the glass-box model's diagram. Points hugging the diagonal
are well calibrated; the companion histogram reveals that the vast majority of windows
sit at low predicted probability --- unsurprising, since freezing is rare --- so the
low-confidence bins are the ones that must be trustworthy, and they are.

\begin{takeaway}
Sigmoid turns scores into probabilities; BCE punishes confident errors; gradient descent
finds the weights; dropout and BatchNorm stop the model memorising; and --- the part that
makes the headline numbers believable --- \textbf{Leave-One-Subject-Out testing measures
the freeze, not the person}. Tested honestly, FoGNet catches $71\%$ of freezes while
leaving $84\%$ of calm gait undisturbed.
\end{takeaway}
